\documentclass[11pt, letterpaper]{article}
\usepackage{afp-style}
\afpdoctype{State Committee Bylaws}

\begin{document}

\afptitlepage{State Committee\\Bylaws Template}{Governing Rules for the State Coordinating Body}

\tableofcontents
\newpage

\section{Name and Authority}

\subsection{Name}
This organization shall be known as the \textbf{Anti-Federalist Party State Committee of \rule{4cm}{0.4pt}} (hereinafter ``the State Committee'').

\subsection{Authority}
The State Committee operates under the authority granted by Article~VI of the national Party Constitution. It serves as the State Coordinating Body referenced therein.

\subsection{Relationship to Chapters}
The State Committee has \textbf{no authority over individual chapters}. Its role is coordination, not control. Chapters retain full autonomy over local operations, endorsements, and governance.

\section{Purpose}

The State Committee exists to:

\begin{enumerate}
  \item Coordinate statewide campaigns and messaging.
  \item Share resources and best practices among chapters.
  \item Represent the state party to the National Council.
  \item File required paperwork with the state election authority.
  \item Maintain compliance with state campaign finance law.
  \item Support chapters with organizing, training, and candidate development.
\end{enumerate}

\section{Membership}

\subsection{Composition}
The State Committee shall consist of:
\begin{enumerate}
  \item One (1) delegate from each chartered chapter within the state.
  \item The elected state officers (Chair, Vice Chair, Secretary, Treasurer).
  \item Up to three (3) at-large members appointed by the State Chair, confirmed by the Committee.
\end{enumerate}

\subsection{Eligibility}
All State Committee members must be members of the Anti-Federalist Party in good standing.

\section{Officers}

\subsection{State Officers}
The State Committee shall elect the following officers:
\begin{enumerate}
  \item \textbf{State Chair} --- Serves as the party's chief representative at the state level, presides over State Committee meetings, and represents the state to the National Council.
  \item \textbf{State Vice Chair} --- Assumes the duties of the Chair in their absence.
  \item \textbf{State Secretary} --- Maintains records, minutes, and correspondence. Files required documents with the state election authority.
  \item \textbf{State Treasurer} --- Manages state party finances, files required campaign finance reports, and maintains transparent records.
\end{enumerate}

\subsection{Terms}
\begin{enumerate}
  \item All state officers serve two-year terms.
  \item No officer may serve more than three consecutive terms in the same position.
  \item Elections shall be held at the annual State Convention.
\end{enumerate}

\subsection{Removal}
Any state officer may be removed by a two-thirds vote of the State Committee, following notice and opportunity to respond.

\section{Meetings}

\subsection{Regular Meetings}
The State Committee shall meet at least quarterly, with at least two meetings conducted virtually.

\subsection{State Convention}
\begin{enumerate}
  \item An annual State Convention shall be held for the purpose of electing officers, adopting resolutions, and coordinating statewide strategy.
  \item Each chapter shall be entitled to at least one delegate.
  \item The convention is the highest authority of the state party.
\end{enumerate}

\subsection{Quorum}
A quorum shall consist of delegates from a majority of chartered chapters or one-third of State Committee members, whichever is greater.

\subsection{Voting}
\begin{enumerate}
  \item Each chapter delegate has one vote.
  \item At-large members have one vote each.
  \item State officers have one vote each, except the Chair, who votes only to break ties.
  \item Ordinary business requires a simple majority. Bylaw amendments require two-thirds.
\end{enumerate}

\section{Candidate Endorsement}

\subsection{State Office Endorsements}
\begin{enumerate}
  \item Endorsements for state legislative or executive office require a two-thirds vote of the State Committee.
  \item All endorsed candidates must affirm the Principles in Article~II of the Party Constitution.
  \item No candidate shall be endorsed who accepts corporate PAC money.
\end{enumerate}

\subsection{Federal Office Endorsements}
Endorsements for federal office (U.S. House, U.S. Senate) require approval by the State Committee and notification of the National Council.

\section{Finances}

\subsection{Revenue}
The State Committee may raise funds through individual donations, chapter contributions, and fundraising events. All revenue restrictions in Article~VIII of the national Constitution apply.

\subsection{Reporting}
\begin{enumerate}
  \item The State Treasurer shall file all reports required by state campaign finance law.
  \item A quarterly financial report shall be presented to the State Committee.
  \item An annual financial summary shall be published to all chapters.
  \item All financial records are open to inspection by any party member.
\end{enumerate}

\subsection{Prohibited Revenue}
The State Committee shall not accept corporate donations, PAC money, or any contribution that creates a quid pro quo obligation.

\section{Amendments}

These bylaws may be amended by a two-thirds vote at the State Convention or by a two-thirds vote of the State Committee, provided that thirty (30) days written notice of the proposed amendment has been given to all members.

\section{Dissolution}

The State Committee may be dissolved by a three-fourths vote of the State Convention. Upon dissolution, assets shall be distributed to the chartered chapters within the state or, in their absence, to the national organization.

\newpage

\afpcertification{State Committee Bylaws}

\vspace{2.5cm}

\signatureline{State Chair}

\signatureline{State Secretary}

\vspace{2cm}

\begin{center}
  {\small\color{afp-gray}\sffamily
    Adopted by the State Convention of \rule{4cm}{0.4pt}\\[0.5em]
    \texttt{anti-federalists.com}
  }
\end{center}

\end{document}
